\section{Discussion and limitations}\label{sec:limits}

\subsection{What the results show}
On maps released in 2026 and within the size range of this benchmark, MICA gives accurate models when it completes: the median TM-score is \TmMedMica{}, and \NTmGeNinetyMica{} of its \NModelMica{} models reach 0.90 or more. With the same AF3 predictions, EModelX(+AF) is indistinguishable from it (mean difference \DiffTmEmodelx{}); E3-CryoFold, which has no AF3 input, is less accurate (\DiffTmEthree{}) and fails more often (\NFailedEthree{} of \NEntries{} entries). ModelAngelo, which also has no AF3 input, matches MICA on the maps up to 3.0~\AA{} (\DiffTmAngeloLow{}) and is lower above 3.0~\AA{} (\DiffTmAngeloHigh{}); part of this gap is an effect of how the TM-score treats fragments, because over all entries the mean difference shrinks from \DiffTmAngelo{} to \DiffFnsAngelo{} with the fragment-agnostic score (Section~\ref{sec:othermeasures}). EMProt, which also uses the AF3 models but builds from the map first, is higher than MICA by a small and consistent margin (\DiffTmEmprot{}); it appeared after MICA, which may explain part of the gap, but this benchmark cannot separate the reasons. After the Holm correction only the difference to EMProt is clear on both tests (Wilcoxon \PwHolmTmEmprot{}, $t$-test \PtHolmTmEmprot{}); E3-CryoFold is lower on the Wilcoxon test (\PwHolmTmEthree{}) but not on the $t$-test (\PtHolmTmEthree{}), and ModelAngelo above 3.0~\AA{} is borderline (\PwHolmTmAngeloHigh{} and \PtHolmTmAngeloHigh{}). The other differences are within the noise of 17 entries.

The step that uses the AF3 structures lifts the median TM-score from \MicaVarInitTmMed{} (the C$\alpha$ model of the network) to \MicaVarAllTmMed{} (Table~\ref{tab:variants}); the AF3 predictions were not scored on their own, so how much of the result is the prior and how much is the map is not separated here. The prior is good: the median over the predictions of the mean pLDDT is \AfPlddtMed{}, and for \NIdentHigh{} of the \NEntries{} entries a chain at 95\,\% identity or more to an older PDB entry exists, which AF3 (whose training data end in 2021; \citealp{abramson2024af3}) may have seen. The benchmark therefore measures how well a method turns a good prior and a new map into a model; the \NIdentLow{} entries with no relative in older PDB entries (9TKA, 9VBD, 26DK) are too few to say how MICA behaves without that prior (MICA failed on one of them; its scores on the others are in Table~\ref{tab:perentry}).

\subsection{Limitations}
\begin{enumerate}
\item \textbf{Sample and range.} There are \NEntries{} entries, \NModelMica{} MICA models and at most \NPairsTmAngelo{} pairs per comparison, all single-particle, protein-only maps of up to \ResiduesMax{} residues and three unique chains (the test maps of the MICA paper go up to several thousand residues). Differences of a few hundredths of TM-score cannot be resolved. Nothing here says how the methods behave on large complexes, where the memory and time limits of the free machine, which already stopped \FailKillMica{} MICA runs for lack of RAM and \FailTimeoutMica{} at its time limit here, would matter even more.
\item \textbf{Docking.} MICA's own docking (Phenix molecular replacement without time limit) could not run on a free machine. The search here is a coarse Fourier search at 6 or 10~\AA{} with a time limit; \DockModePartial{} entries are partly docked and \DockModeNone{} not at all, and the placements were not checked against the deposited models. The indirect evidence (Section~\ref{sec:dockresults}) suggests that this understates MICA, but it rests on \DockEffectN{} entries and on a descriptive split.
\item \textbf{No real-space refinement.} MICA's last step stopped with a parameter error in Phenix~\PhenixVersion{}, for every entry. The scored models are MICA's models before refinement. I did not refine the outputs of the other methods either, but their own pipelines may include refinement-like steps, so the comparison is not strictly symmetric; the absolute numbers are not those of the authors' full pipeline, and the effect on all-atom quality was not measured.
\item \textbf{Memory patch and reruns.} The network of MICA predicts every window on its own, as MICA does for maps with 200 windows or fewer, so that the larger maps fit into the 15~GB of the GPU; MICA's own setting (batches of eight windows) was not compared with it, and the setting is meant to change memory use only. The \NRerunMica{} models of \RerunMicaList{} come from reruns after a first failure with MICA's default batching, an exception to the one-run rule that is reported.
\item \textbf{Failures depend on the machine.} Of the \FailTotal{} runs that ended without a model, \FailKill{} were kills by the system (probably for lack of RAM), \FailTimeout{} reached a run limit (\FailTimeoutMica{} of MICA and \FailTimeoutAngelo{} of ModelAngelo) and \FailCrashEthree{} were crashes of E3-CryoFold (exit code 1; entries \FailCrashListEthree{}). All happened on a machine with 31~GB of RAM, a 15~GB GPU and fixed run limits; a larger machine would probably have completed the kills and perhaps the runs that reached a limit, but not the crashes. Scoring failures as zero therefore describes this machine and the programs together.
\item \textbf{One run per entry.} Each method ran once per entry, with one AF3 prediction (one random seed, server defaults) per chain; the variation between AF3 seeds and between repeated runs was not measured. Run times come from shared free machines and are indicative.
\item \textbf{Reference models and measures.} The deposited models are not ground truth: they contain errors and may have been built with tools such as ModelAngelo or from AF predictions, which could favour similar methods. The maps often hold more copies than the deposit, so methods that build every copy have a low C$\alpha$ match by construction; US-align's TM-score counts one fragment per chain. The conclusions rest on the TM-score and the C$\alpha$ quality score and are checked with the fragment-agnostic score.
\item \textbf{Protocol choices.} EModelX(+AF) was run with the AF3 predictions as templates, not in its default AlphaFold mode, and E3-CryoFold only in its \texttt{pre\_align} protocol; ModelAngelo had no AlphaFold input by design. These choices were fixed before the runs. The newer methods are two, because CryoAtom could not run within the compute budget, and the literature search could not use forward citations of MICA.
\item \textbf{Selection.} The map--model correlation rule (CC$_\text{mask}$ of at least \SelCcMin{}) removed \SelCcRejected{} entries, all in the stratum above 3.0~\AA{}, and \SelCcFailed{} more could not be computed; the stratum above 3.0~\AA{} therefore holds maps that are well explained by their models. The seeded order favours common specimens to the degree that the pool does.
\end{enumerate}
